% General theory: partially identified global portraits and downstream decisions.

The mapping calculus above can be stated more generally. Scientific sources
rarely reveal a unique global state. They reveal projections of that state,
with different coordinates observed, suppressed, coarsened, or measured under
different operationalizations. The appropriate integration object is
therefore not always one completed graph. It is an envelope of the global
portraits still compatible with the observations and their source-bound
constraints. This section develops that object and separates three questions:
what the sources identify, what additional information could identify, and
what a downstream decision may safely do before identification is complete.

The construction imports the logic of partial identification
\cite{manski2003partial,tamer2010partial}, comparison of information
\cite{blackwell1953equivalent}, and robust decision theory
\cite{berger1985decision,walley1991imprecise}. The contribution is not a new
claim to those donor theories. It is their claim-relative composition with
provenance-bound scientific projections, typed mapping constraints, and
global-portrait construction.

\subsection{Worlds, observation fibres, and portrait envelopes}

Let $\Omega$ be a set of scientifically admissible worlds. A world
$\omega\in\Omega$ fixes the source records, the source-to-projection
relations, and a global portrait $G(\omega)\in\mathcal G$. The observation
operator
\begin{equation}
  O:\Omega\longrightarrow\mathcal Y
\end{equation}
returns exactly the evidence available to the integration procedure: the
observed projection coordinates, explicit missingness, whatever source
identity/version information the interface exposes, and the mapping
constraints licensed by those sources. It does not fill an unobserved
coordinate from the world that generated it.

The operator is interface-relative. If a procedure receives source text and
semantic provenance, those are part of $O$; if it receives only projections,
they are not. Opaque record identifiers that carry no licensed scientific
meaning are quotiented by bijective relabelling rather than treated as a
lookup key for memorising gold. Every impossibility statement must name this
observation boundary.

\begin{definition}[Observation fibre and epistemic portrait envelope]
For observed evidence $y\in\mathcal Y$, define
\begin{equation}
  \Omega_y=\{\omega\in\Omega:O(\omega)=y\},\qquad
  \mathfrak E(y)=\{G(\omega):\omega\in\Omega_y\}.
\end{equation}
$\Omega_y$ is the observation fibre and $\mathfrak E(y)$ is the epistemic
portrait envelope. A query $q:\mathcal G\to\Theta$ has identified set
\begin{equation}
  \Theta_q(y)=\{q(g):g\in\mathfrak E(y)\}.
\end{equation}
The query is point identified at $y$ exactly when $\Theta_q(y)$ is a
singleton.
\end{definition}

Write $\mathcal Y_{\mathrm r}=O(\Omega)$ for the set of realised
observations. If $y\notin\mathcal Y_{\mathrm r}$, then $\Omega_y$ and
$\Theta_q(y)$ are empty. We treat such an input as an explicit
\textsc{invalid model/interface} terminal, denoted $\bot_{\mathrm{inv}}$,
rather than as an ordinary scientific envelope. In particular, no robust loss
is optimized by taking a supremum over the empty set. All decision statements
below are restricted to $y\in\mathcal Y_{\mathrm r}$ unless an explicit
invalid terminal is mentioned.

The fibre calculus is set-theoretic and pointwise. A globally measurable
decoder is an additional conclusion, not a consequence of fibre constancy on
arbitrary measurable spaces. If $\Omega,\mathcal Y$, and the query codomain
carry sigma-algebras, such a conclusion requires a measurable factorisation
$q\circ G=d\circ O$. Where conditional probabilities are invoked below, they
are separately licensed probabilistic objects. If derived from a joint law,
we assume standard-Borel spaces and a declared version of the regular
conditional law. The set-valued portrait envelope alone does not manufacture
a conditional probability.

The query can ask whether two projections may be glued, which typed
obstruction holds, what numerical estimand a merged record implies, or which
downstream action is warranted. Identification is therefore claim-relative:
an envelope can identify one query while leaving another unresolved. A
portrait need not be completed in every coordinate before it licenses a
specific conclusion.

The envelope also separates scientific plurality from epistemic uncertainty.
Plurality is a feature represented \emph{inside} a feasible portrait: several
incompatible but source-valid views coexist. Epistemic uncertainty is
variation \emph{across} feasible portraits because the available observation
does not decide which completion obtains. A plural-view terminal can therefore
be point identified; an unresolved terminal cannot.

\subsection{The fibre criterion and an impossibility result}

\begin{theorem}[Fibre criterion for scientific identification]
\label{thm:fibre}
For any observation $y$ with $\Omega_y\neq\varnothing$ and query $q$, the
following are equivalent:
\begin{enumerate}
  \item $q$ is point identified at $y$;
  \item $q\circ G$ is constant on $\Omega_y$;
  \item there exists an observation-only answer function $d(y)$ that equals
        $q(G(\omega))$ for every $\omega\in\Omega_y$.
\end{enumerate}
If two worlds in the same observation fibre have different query values, no
deterministic procedure reading only $y$ can be correct in both worlds.
\end{theorem}

\begin{proof}
The image of $\Omega_y$ under $q\circ G$ is exactly $\Theta_q(y)$. It is a
singleton if and only if the map is constant on the fibre. In that case its
unique value defines $d(y)$. Conversely, if $d(y)$ is correct for every world
in the fibre, all those worlds must have query value $d(y)$. The final claim
follows immediately: the procedure receives the same input in both worlds and
therefore emits the same answer, which cannot equal two different values.
\end{proof}

This theorem turns a partial-observation failure into a scientific result
rather than an invitation to invent a fifth rule over the same information.
When a coordinate erased by extraction is the only fact distinguishing two
source relations, model capacity, prompt changes, and a more elaborate
comparison rule cannot restore it. The remedy must either acquire information,
declare an assumption that removes worlds from the fibre, or return a
set-valued/unresolved answer.

\subsection{A canonical sufficient interface}

The envelope is not a demand that every integration system share one internal
architecture. It defines the least claim-relative interface a downstream
consumer must be able to recover.

\begin{definition}[Envelope sufficiency]
A summary $S:\mathcal Y\to\mathcal Z$ is sufficient for query $q$ when there
exists a decoder $h$ such that
\begin{equation}
  \Theta_q(y)=h(S(y))\qquad\text{for every }y\in\mathcal Y_{\mathrm r}.
\end{equation}
Thus equal summaries may never hide different identified sets.
\end{definition}

\begin{theorem}[Universal factorisation of sufficient interfaces]
\label{thm:factorisation}
On $\mathcal Y_{\mathrm r}$, the map $T_q:y\mapsto\Theta_q(y)$ is sufficient
for $q$. Every other sufficient summary $S$ refines it: there is a map $h$
with $T_q=h\circ S$. Consequently, all sufficient implementations evaluated
only through the same set-robust loss have the same robust loss table and the
same set of minimax actions, up to ties.
\end{theorem}

\begin{proof}
$T_q$ is decoded by the identity map on realised observations. The
factorisation for another sufficient $S$ is its defining decoder. Robust upper
loss for action $a$ is a function only of the recovered set,
$\sup_{\theta\in T_q(y)}L(a,\theta)$. Equal recovered sets therefore give the
same value for every action, and minimisation gives the same optimal-action
set.
\end{proof}

This theorem explains both the ambition and the limit of the architecture. A
semantic product that exposes exactly the same identified set is not an
inferior baseline; it is an information-equivalent implementation and must
tie at this interface. A product that collapses two observations with
different identified sets is insufficient for that query, no matter how rich
its remaining representation is. Any empirical advantage must therefore come
from a genuinely finer observation, justified constraints, a more faithful
approximation to the envelope, or a different declared downstream loss---not
from the name or central location of the layer.
The conclusion is deliberately about the displayed set-robust rule. Two
summaries that recover the same identified set may still support different
Bayes decisions if one additionally carries a licensed probability law over
that set.

\subsection{Information refinement and the observation floor}

Say that $O_2$ refines $O_1$ when there is a map $r$ such that
$O_1=r\circ O_2$. Thus $O_2$ retains at least as much information as $O_1$;
the relation is deterministic and makes no probabilistic sufficiency claim.
Let $\mathcal A$ be a finite nonempty downstream action set and
$L:\mathcal A\times\Theta\to[0,\infty)$ a declared loss. Define the robust
decision floor, for $y\in\mathcal Y_{\mathrm r}$,
\begin{equation}
  B_q(y)=\min_{a\in\mathcal A}\ \sup_{\theta\in\Theta_q(y)}L(a,\theta).
  \label{eq:decision-floor}
\end{equation}

\begin{theorem}[Information-refinement monotonicity]
\label{thm:refinement}
If $O_2$ refines $O_1$, then for every realised world $\omega$,
\begin{equation}
 \Theta^{(2)}_q(O_2(\omega))
 \subseteq
 \Theta^{(1)}_q(O_1(\omega))
 \quad\text{and}\quad
 B^{(2)}_q(O_2(\omega))\leq B^{(1)}_q(O_1(\omega)).
\end{equation}
\end{theorem}

\begin{proof}
Every world having the refined observation $O_2(\omega)$ also has coarse
observation $r(O_2(\omega))=O_1(\omega)$, so the refined fibre is a subset of
the coarse fibre. Taking images under $q\circ G$ preserves inclusion. For any
fixed action, the supremum of its loss over a subset cannot increase; taking
the minimum over actions preserves the inequality.
\end{proof}

The quantity $B_q(y)$ is an observation floor, not a defect score for a
particular algorithm. It can fall when truthful extraction, measurement,
provenance, or source acquisition refines the observation. It cannot be made
to disappear merely by changing the decision rule over the same fibre. For a
candidate rule $\delta$, its floor-adjusted avoidable harm is
\begin{equation}
 X_\delta(y)=
 \sup_{\theta\in\Theta_q(y)}L(\delta(y),\theta)-B_q(y)\geq 0.
 \label{eq:excess-harm}
\end{equation}
This is the quantity a successor experiment should compare. Counting total
error without subtracting or stratifying by the observation floor confounds
avoidable decision failure with information that no compared rule received.

\subsection{Coordinate opportunity is within-comparison support}

The same fibre logic applies to experimental support. Suppose a protected
evaluation is organised into comparison blocks $i=1,\ldots,n$, with two
source-grounded members $y_i^L,y_i^R$ in each block. For scientific
coordinate $c$, let $z_c(y)$ be its candidate-visible value and define the
coordinate-opportunity matrix
\begin{equation}
 S_{ic}=\mathbf 1\!\left\{z_c(y_i^L)\neq z_c(y_i^R)\right\}.
 \label{eq:coordinate-opportunity}
\end{equation}
This is a support condition, not an outcome. A coordinate may take many
different values over the atlas while its entire column of $S$ is zero: the
value changes between blocks but never supplies a comparison inside one.

\begin{theorem}[No coordinate attribution without comparison support]
\label{thm:coordinate-support}
Fix a coordinate $c$ and a paired design $\mathcal D=\{(y_i^L,y_i^R)\}_{i=1}^n$.
If $S_{ic}=0$ for every $i$, then any two decision rules that agree on the
observed design support are observationally indistinguishable. In particular,
when the admissible domain contains an off-support pair with a $c$-contrast,
the design cannot distinguish a rule that is insensitive to $c$ from a rule
that differs only on such pairs. Consequently, marginal variation of $z_c$ across blocks does
not identify coordinate-specific decision value, harm reduction, or
necessity without an additional cross-block exchangeability or structural
extrapolation assumption.
\end{theorem}

\begin{proof}
The first claim follows because a statistic of the observed blocks depends
only on rule values on that support. For the particular construction, let
$\delta_0$ be any rule and let $\delta_1$ equal it on every observed pair but
differ on one admissible off-support pair with a $c$-contrast. The hypothesis
$S_{ic}=0$ ensures that this difference is never queried by $\mathcal D$, so
the two rules induce exactly the same outputs and losses on every observed
block. No statistic of those blocks can distinguish them.
Attributing a difference to $c$ from between-block variation therefore
requires an assumption linking otherwise different blocks; that assumption
is not supplied by the paired design itself.
\end{proof}

The theorem separates three increasingly strong gates. A field being
populated is weaker than marginal variation; marginal variation is weaker
than nonzero within-block contrast; and within-block contrast is weaker than
coordinate-specific identification when several coordinates move together.
For a claim about coordinate $c$, the protocol must therefore report the
column sum $\sum_i S_{ic}$, the joint pattern of the other coordinate columns,
and the protected query or loss opportunity. A zero column is an automatic
\textsc{cannot check} for that coordinate. A nonzero column admits analysis
but does not by itself prove a causal coordinate effect.

\subsection{Credal risk envelopes for merge, split, and unresolved actions}

For a realised observation $y$ at which probabilistic authority is available,
let the query space be measurable, let $\Theta_q(y)$ be measurable, and let
$\mathcal K_y$ be a nonempty set of licensed probability measures satisfying
$P(\Theta_q(y))=1$ for every $P\in\mathcal K_y$. For each
$a\in\mathcal A$, assume $L(a,\cdot)$ is measurable and define its lower and
upper expected risks
\begin{equation}
 \underline R(a\mid y)=
 \inf_{P\in\mathcal K_y}\int L(a,\theta)\,dP(\theta),
 \qquad
 \overline R(a\mid y)=
 \sup_{P\in\mathcal K_y}\int L(a,\theta)\,dP(\theta).
 \label{eq:credal-risk-envelope}
\end{equation}
These values may be extended-real for a general nonnegative loss. A robust
rule minimizes $\overline R(a\mid y)$. Because $\mathcal A$ is finite and
nonempty, the pointwise minimum is attained. A globally measurable robust
rule additionally requires the upper-risk functions to be measurable; with a
fixed ordering of $\mathcal A$, measurable upper risks give a measurable
tie-broken selector.

This credal rule contains the earlier set-robust floor as a special case. If
$\mathcal K_y$ is the class of all probability measures supported on
$\Theta_q(y)$, its upper expected risk equals
$\sup_{\theta\in\Theta_q(y)}L(a,\theta)$ because all point masses are
available. More generally, replacing $\mathcal K_y$ by a nonempty subset can
only decrease every upper risk and the resulting minimax floor. This is a
refinement of licensed probabilistic authority, distinct from refining the
observation operator itself.

For a binary mapping query, let $\theta=1$ mean that glue is licensed and
$\theta=0$ mean that it is not. The downstream actions are merge $M$, keep
separate $S$, and unresolved $U$. A false merge costs $c_{\mathrm{FM}}\geq0$, a
false split costs $c_{\mathrm{FS}}\geq0$, and deferral costs $c_U\geq0$. Because
a binary law is determined by $p=P(\theta=1)$, write
\begin{equation}
 K_y=\{P(\theta=1):P\in\mathcal K_y\}\subseteq[0,1],
 \qquad \ell_y=\inf K_y,\quad u_y=\sup K_y.
 \label{eq:binary-credal-set}
\end{equation}
The set $K_y$ need not be an interval, convex, closed, or endpoint-attaining.

\begin{proposition}[Binary robust action over an arbitrary credal set]
\label{prop:binary-bound}
The closed interval hulls determined by the lower and upper risks are
\begin{align}
 [\underline R(M\mid y),\overline R(M\mid y)]
   &=[c_{\mathrm{FM}}(1-u_y),c_{\mathrm{FM}}(1-\ell_y)],\\
 [\underline R(S\mid y),\overline R(S\mid y)]
   &=[c_{\mathrm{FS}}\ell_y,c_{\mathrm{FS}}u_y],\\
 [\underline R(U\mid y),\overline R(U\mid y)]
   &=[c_U,c_U].
 \label{eq:binary-risk-envelope}
\end{align}
In particular, the worst-case expected losses are
\begin{equation}
 \overline R(M\mid y)=c_{\mathrm{FM}}(1-\ell_y),\quad
 \overline R(S\mid y)=c_{\mathrm{FS}}u_y,\quad
 \overline R(U\mid y)=c_U.
 \label{eq:binary-risk}
\end{equation}
Consequently, unresolved is minimax-optimal exactly when
\begin{equation}
 c_U\leq
 \min\{c_{\mathrm{FM}}(1-\ell_y),c_{\mathrm{FS}}u_y\}.
 \label{eq:defer-bound}
\end{equation}
\end{proposition}

\begin{proof}
Merging has risk $c_{\mathrm{FM}}(1-p)$, a decreasing affine function of $p$.
Its infimum and supremum over $K_y$ are therefore its limiting values at $u_y$
and $\ell_y$, respectively, whether or not those endpoints are attained.
Keeping separate has risk $c_{\mathrm{FS}}p$, so the order is reversed.
Deferral has constant risk. Comparing the three upper risks gives the exact
criterion.
\end{proof}

\begin{corollary}[Sharp interval and vacuous special cases]
If the licensed credal set is exactly the full sharp interval
$K_y=[\underline p(y),\overline p(y)]$, then
\begin{equation}
 \overline R(M\mid y)=c_{\mathrm{FM}}(1-\underline p),\qquad
 \overline R(S\mid y)=c_{\mathrm{FS}}\overline p,
\end{equation}
and Equation~\eqref{eq:defer-bound} holds with
$(\ell_y,u_y)=(\underline p,\overline p)$. With no probabilistic authority
beyond the vacuous set $K_y=[0,1]$ when both binary states remain feasible,
unresolved is minimax-optimal exactly when
$c_U\leq\min\{c_{\mathrm{FM}},c_{\mathrm{FS}}\}$.
\end{corollary}

If $[\underline p,\overline p]$ is merely an outer bound satisfying
$K_y\subseteq[\underline p,\overline p]$, then
$c_{\mathrm{FM}}(1-\underline p)$ and
$c_{\mathrm{FS}}\overline p$ are conservative upper bounds on the two robust
risks. Optimizing against the whole outer interval is a defensible, more
conservative decision problem, and the interval criterion is exact for that
enlarged problem. It does \emph{not} by itself certify that unresolved is
minimax for the unknown smaller $K_y$: the alternative actions' true robust
risks may be strictly below their outer upper bounds. Exact minimax claims
therefore require the actual extrema $(\ell_y,u_y)$ or a separately justified
sharp interval. For example, with $K_y=\{1/2\}$, outer interval $[0,1]$,
$c_{\mathrm{FM}}=c_{\mathrm{FS}}=1$, and $c_U=3/4$, the outer problem selects
unresolved, whereas both determinate actions have true robust risk $1/2$.

The costs are downstream and use-specific. A clinical synthesis may assign a
larger cost to a false merge than a discovery interface; a screening system
may assign a larger cost to a false split. The portrait mechanism should
expose the identified set and provenance needed to evaluate those choices,
not silently encode one universal utility function.

\subsection{Global composition is a joint-completion problem}

Let each accepted source or mapping record impose a set
$C_j\subseteq\Omega$ of worlds consistent with its licensed conditions. The
jointly feasible set is
\begin{equation}
  F(y)=O^{-1}(y)\cap\bigcap_{j=1}^{m}C_j.
  \label{eq:joint-completion}
\end{equation}
A global glue requires $F(y)\neq\varnothing$; a claim additionally requires
its query to be constant on $F(y)$. Pairwise mapping success does not imply
either condition.

\begin{proposition}[Pairwise compatibility is insufficient]
\label{prop:pairwise}
There exist three mapping constraints that are pairwise jointly satisfiable
but have no common global completion. Hence a system that accepts every
pairwise-compatible edge can license an inconsistent global portrait.
\end{proposition}

\begin{proof}
Take $\Omega=\{1,2,3\}$ and
$C_1=\{1,2\}$, $C_2=\{2,3\}$, $C_3=\{1,3\}$. Every pairwise intersection is
nonempty, but $C_1\cap C_2\cap C_3=\varnothing$.
\end{proof}

Constraint processing and acyclic join theory provide stronger algorithms for
deciding global consistency than this elementary counterexample
\cite{dechter2003constraint}. The point here is architectural: pairwise
ontology or schema correspondences are proposals for $C_j$, not self-executing
authority for a global merge. Empty joint completion is an identified global
obstruction. Multiple completions with different query values are partial
identification. Multiple source-valid views inside every completion are an
identified plural portrait.

\subsection{Positive-reference absence without closure authority}

Many discovery-style evaluations publish a set of asserted positives rather
than a complete binary labelling of a candidate universe. Let $U$ be a frozen
candidate universe and let $R^+\subseteq U$ be the observed positive reference.
A binary completion is a map
$t:U\rightarrow\{\mathsf{glue},\mathsf{obstruction}\}$ satisfying
$t(x)=\mathsf{glue}$ for every $x\in R^+$. The licensed completion class may
also contain source-grounded negative assertions or an authority constraint
certifying that $R^+$ is exhaustive. Absent such a constraint, omission is an
observation, not a negative label.

\begin{theorem}[Positive-reference non-identification without closure authority]
\label{thm:positive-reference-nonidentification}
Fix $x\in U\setminus R^+$. Suppose the licensed completion class retains both
(i) a completion in which $t(x)=\mathsf{glue}$ and the positive reference
omitted that true equivalence and (ii) a completion in which
$t(x)=\mathsf{obstruction}$. Then the binary truth at $x$ is not identified by
$(U,R^+)$: its identified set is
\begin{equation}
  I_{R^+}(x)=\{\mathsf{glue},\mathsf{obstruction}\}.
\end{equation}
Consequently, reference absence alone cannot license an obstruction label or
an unconditional point-valued harm score. Assigning obstruction to every
$x\notin R^+$ adds a closed-world exhaustivity assumption.
\end{theorem}

\begin{proof}
The two licensed completions induce the same observed pair $(U,R^+)$ but give
different binary answers at $x$. The fibre criterion therefore makes both
answers members of the identified set and rules out any uniformly correct
singleton decoder based only on $(U,R^+)$. Declaring the second answer unique
removes the first completion; that removal is exactly additional negative or
exhaustivity authority, not a consequence of reference omission.
\end{proof}

The premise is deliberately authority-relative. The theorem does not assert
that every reference alignment is incomplete, nor does a missing licence file
prove incompleteness. It states what follows when the admissible-world model has
not licensed closure and therefore retains both completions. Explicit
disjointness evidence, source-certified negatives, logical constraints, or an
authoritative exhaustivity certificate over the frozen universe may remove one
completion. Matcher nonselection and agreement among correlated algorithms do
not do so by themselves.

This result is not specific to ontology matching. It applies whenever a
benchmark records discoveries without certifying the full negative complement.
In that setting the appropriate target is an authority-aware identification
envelope: point evaluation may be reported only for identified cells, while
unlicensed omissions remain set-valued or \textsc{CannotCheck}.

\subsection{Direct-certificate calibration}

Let $C_G,C_O\subseteq U$ be licensed direct-certificate sets for glue and
obstruction.  Consider the completion class of all binary maps
$t:U\rightarrow\{\mathsf{glue},\mathsf{obstruction}\}$ satisfying
$t(x)=\mathsf{glue}$ on $C_G$ and $t(x)=\mathsf{obstruction}$ on $C_O$, with
no negative label inferred from certificate absence.

\begin{corollary}[Direct-certificate identified sets]
\label{cor:direct-certificate-identification}
If $C_G\cap C_O=\varnothing$, the sharp identified set is
\begin{equation}
 I_C(x)=
 \begin{cases}
  \{\mathsf{glue}\}, & x\in C_G,\\
  \{\mathsf{obstruction}\}, & x\in C_O,\\
  \{\mathsf{glue},\mathsf{obstruction}\}, & x\notin C_G\cup C_O.
 \end{cases}
\end{equation}
Any conflict-free addition of direct certificates can only shrink these sets.
If $x$ remains uncertified, absence alone cannot point-identify obstruction.
If $x\in C_G\cap C_O$, the direct authority record is contradictory and must
be retained as \textsc{CannotCheckConflict} rather than silently resolved.
\end{corollary}

\begin{proof}
On a singly certified pair, every licensed completion has the certified value.
On an uncertified pair, changing only $t(x)$ constructs licensed glue and
obstruction completions with the same certificate observation, so the fibre
criterion retains both values. Adding a nonconflicting certificate removes
completions but adds none, hence identified sets shrink monotonically.
\end{proof}

This is a direct specialization of the general fibre and identified-set
results, not a claim that direct ontology axioms exhaust semantic truth. Its
role is to expose a finite authority domain on which binary scoring is licensed
without converting nonselection into a negative label.

\subsection{The sharp binary covering envelope}

For a retained observation $x$, let
$F_x=\{\omega:O(\omega)=x\}$ and let
$I(x)=\{q(\omega):\omega\in F_x\}\subseteq
\{\mathsf{glue},\mathsf{obstruction}\}$ be the binary identified set. An
envelope rule $E$ has distribution-free coverage on the licensed world class
when $q(\omega)\in E(O(\omega))$ for every licensed world.

\begin{corollary}[Minimal binary covering envelope]
\label{cor:binary-covering-envelope}
An envelope rule has distribution-free coverage if and only if
$I(x)\subseteq E(x)$ for every attained observation $x$. Consequently $I$ is
the unique inclusion-minimal covering envelope. If one observation fibre
contains both a glue world and an obstruction world, no singleton action is
coverage-valid without an additional identification assumption.
\end{corollary}

\begin{proof}
If $E$ covers every licensed world and $y\in I(x)$, some
$\omega\in F_x$ has $q(\omega)=y$; coverage gives
$y\in E(O(\omega))=E(x)$. The converse follows because
$q(\omega)\in I(O(\omega))\subseteq E(O(\omega))$ for every world.
Inclusion-minimality is immediate.
\end{proof}

This is the standard identified-set argument specialized to the portrait
decision, not a new general theorem. It makes the development boundary exact:
agreement between correlated algorithms cannot by itself certify fibre
constancy, while the full binary envelope attains coverage one by construction
but may impose severe downstream deferral harm. A scientifically useful
successor must license a proper subenvelope from additional, independently
validated structure while retaining coverage.

\subsection{Evaluation implied by the theory}

The envelope theory changes the scientific target from forced-label accuracy
to four jointly reported quantities:
\begin{enumerate}
  \item \emph{validity}: whether protected source-grounded truth is excluded
        from the returned envelope;
  \item \emph{sharpness}: the size or decision diameter of the envelope,
        reported only together with validity;
  \item \emph{floor-adjusted decision harm}: Equation~\eqref{eq:excess-harm}
        under a declared downstream loss; and
  \item \emph{information value}: the reduction in the identified set or
        decision floor from a specified observation refinement.
\end{enumerate}
These targets envelop the original glue-or-obstruction calculus. A singleton
envelope recovers an ordinary determinate mapping; a non-singleton envelope
explains why unresolved is required; an empty joint-completion set certifies a
global obstruction. The current experiments cover only special cases of this
theory. They do not by themselves establish calibrated naturalistic envelopes
or downstream utility.

\subsection{Authority-augmented identification and composition}
\label{sec:authority-augmented}

The portrait envelope above identifies content relative to an observation.
Some scientific queries are additionally undefined unless the evidence
possesses claim-relevant authority: a source and version are bound, the
material may lawfully support the declared use, the method and comparator can
execute, the population and semantics coincide, and the outcome is evaluated
under the declared custody. Missing such a condition is not a third truth
value and is not evidence against the method. It is an authority terminal.

For a query $q:\Omega\to\Theta$, freeze the finite set $J_q$ of authority
gates required by the claim and let $a_j:\Omega\to\{0,1\}$ be their truth in a
world. Introduce $\bot_{\mathrm{auth}}\notin\Theta$ and define
\begin{equation}
 q^A(\omega)=
 \begin{cases}
 q(\omega),&\prod_{j\in J_q}a_j(\omega)=1,\\
 \bot_{\mathrm{auth}},&\text{otherwise},
 \end{cases}
 \qquad
 \Theta^A_q(y)=\{q^A(\omega):O(\omega)=y\}.
 \label{eq:authority-envelope}
\end{equation}
The terminal is the formal counterpart of \textsc{CannotCheck}. In
particular, if all compatible completions lack a required executable or
reference identity, then $\Theta^A_q(y)=\{\bot_{\mathrm{auth}}\}$ rather than
the empty set or a negative performance value.

\begin{theorem}[Authority-aware fibre criterion]
\label{thm:authority-fibre}
For a realised observation $y$ and $\theta\in\Theta$, the authorized claim
value is point identified as $\theta$ if and only if
$\Theta^A_q(y)=\{\theta\}$. Equivalently, every authority gate is true and
$q(\omega)=\theta$ throughout the observation fibre. If
$\bot_{\mathrm{auth}}\in\Theta^A_q(y)$, the observation does not identify the
authority needed for the claim; if two ordinary values occur, it does not
identify the claim's substance.
\end{theorem}

\begin{proof}
$\Theta^A_q(y)$ is the image of the observation fibre under $q^A$. It is the
ordinary singleton $\{\theta\}$ exactly when $q^A$ is constantly $\theta$ on
that fibre. By Equation~\eqref{eq:authority-envelope}, this holds exactly when
all frozen gates equal one and $q$ is constantly $\theta$ there. The two
failure cases follow from whether the image contains the authority terminal or
multiple ordinary values.
\end{proof}

The same refinement order applies. If $O_2$ refines $O_1$ while $q$ and the
frozen authority predicates are unchanged, the refined fibre is a subset of
the coarse fibre and therefore
\begin{equation}
 \Theta_q^{A,(2)}(O_2(\omega))
 \subseteq
 \Theta_q^{A,(1)}(O_1(\omega)).
 \label{eq:authority-refinement}
\end{equation}
Truthful acquisition of a rights fact, executable revision or custodied
outcome can therefore shrink the set. A mismatched repository, changed
population or outcome-informed adapter is a changed scientific identity, not
a refinement of the frozen observation.

Local readiness is not sufficient for global composition. Consider two worlds
with the same observed stage artifacts and with every local source, rights,
parser and entrypoint gate true. Let an unobserved relation say whether the two
artifacts belong to the same scientific identity, with value zero in one world
and one in the other, and let the end-to-end query be that relation. Every
local statement is fixed and authorized, but the global identified set is
$\{0,1\}$. A passed collection of local preflights therefore yields an
end-to-end result only when the global query factors through those outputs
under an explicit composition relation.

For comparator superiority, the required gates include both executable
identities, resource and configuration completeness, use rights, task and
population identity, reference authority, scorer custody and terminal-
preserving completion. If any is absent or unresolved, a missing run cannot be
imputed as a loss, tie, abstention, obstruction or zero. Native artifact smoke
may identify a narrower executability query, but it does not identify
superiority without the remaining truth, population and custody gates. This
extends the content envelope to scientific evaluation without claiming a new
general theory beyond the fibre criterion.

\endinput
